\documentclass{article}
\usepackage{/globals/math}
\usepackage{/globals/diagrams}
\usepackage{/globals/anki}
\ankiprefix{transistoren}
\ankitopic{Transistoren}
\ankishow
\begin{document}
\section{Transistoren} 
Computerhardware besteht aus Transistoren; elektrischen Schaltern die entweder an oder aus sein können.

Transistoren haben eine Quelle (Source), ein Gatter (Gate) und eine Senke (Drain). Ob die Quelle mit der Senke verbunden ist hängt davon ab ob Spannung am Gatter anliegt; bei nMOS Transistoren fließt Strom wenn Spannung anliegt, bei pMOS Transistoren fließt kein Strom wenn Spannung anliegt. Dass ein Transistor ein pMOS Transistor ist wird durch einen unausgefüllten Kreis dargestellt.
\ankinote{bestandteile}{Bestandteile}{Quelle (Source), Gatter (Gate), Senke (Drain)}
\ankinote{funktion}{Funktionsweise}{Quelle und Senke verbunden, abhängig von Gatter}
\ankinote{nmos}{nMOS}{Strom wenn Spannung}
\ankinote{pmos}{pMOS}{Strom wenn keine Spannung, Kreis}

\subsection{Eine Not-Schaltung}
Aus zwei Transistoren kann eine Not-Schaltung gebaut werden.
\ztikz{
 \pmos(0,1)
 \nmos(0,-1)
 \conn(-1, 1)(-1, -1)
 \conn(-1, 0)(-2, 0)
 \conn(0, 2)(0, 1.5)
 \conn(0, -2)(0, -1.5)
 \conn(0, 0.5)(0, -0.5)
 \conn(0, 0)(1, 0)
 \conn(-0.5, 2)(0.5, 2)
 \zlabel(0, 2){above}{$+5 \text{V}$}
 \inters(0, 2)
 \conn(-0.3, -2)(0.3, -2)
 \conn(-0.2, -2.1)(0.2, -2.1)
 \conn(-0.1, -2.2)(0.1, -2.2)
 \zlabel(0, -2.2){below}{Masse}
 \zlabel(-2, 0){left}{$x$}
 \zlabel(1, 0){right}{$\overline{x}$}
}
Wenn an $x$ Spannung anliegt, so wird der obige pMOS blockiert und es besteht keine Verbindung zwischen dem Strom und dem Ausgang, aber der Ausgang und die Masse ist verbunden. Liegt bei $x$ kein Strom an, so ist die Verbindung zur Masse blockiert aber die Verbindung zum obigen Strom frei.

So können durch Verbindungen von Transistoren auch andere Gatter gabaut werden.


\end{document}